% =====================================================================
%  Eclipse 2025.4 — Faults & Non-Neighbour Connections (NNCs)
%  1-to-1 Mathematical Model
%
%  Every equation maps to a numbered equation in:
%    EclipseTechnicalDescription.pdf  (TD chapter "Transmissibility calculations",
%                                      pp. 33-51)
%  or to the RM keyword pages for NNC, EDITNNC, EDITNNCR, MULTFLT.
%
%  Compile with:  pdflatex math-model.tex
%
%  Author: AXON, built 2026-05-23
% =====================================================================

\documentclass[11pt,a4paper]{article}

\usepackage[utf8]{inputenc}
\usepackage[T1]{fontenc}
\usepackage{lmodern}
\usepackage{microtype}
\usepackage{amsmath,amssymb,amsthm,mathtools}
\usepackage{bm}
\usepackage[margin=2.4cm]{geometry}
\usepackage{booktabs}
\usepackage{xcolor}
\usepackage{hyperref}
\hypersetup{colorlinks=true, linkcolor=blue!50!black, urlcolor=blue!60!black}
\usepackage[font=small,labelfont=bf]{caption}
\usepackage{parskip}

\newcommand{\TDref}[1]{\textit{TD Eq.\,#1}}
\newcommand{\RMref}[1]{\textit{RM Eq.\,#1}}
\newcommand{\kw}[1]{\texttt{#1}}
\newcommand{\note}[1]{\par\smallskip\textit{#1}\par\smallskip}

\DeclareMathOperator{\CDARCY}{CDARCY}
\DeclareMathOperator{\TRANX}{TRANX}
\DeclareMathOperator{\TRANY}{TRANY}
\DeclareMathOperator{\TRANZ}{TRANZ}
\DeclareMathOperator{\TRANR}{TRANR}
\DeclareMathOperator{\TMLT}{TMLT}
\DeclareMathOperator{\DIPC}{DIPC}
\DeclareMathOperator{\RNTG}{RNTG}
\DeclareMathOperator{\MULTFLT}{MULTFLT}
\DeclareMathOperator{\PERMX}{PERMX}
\DeclareMathOperator{\PERMY}{PERMY}
\DeclareMathOperator{\PERMZ}{PERMZ}
\DeclareMathOperator{\PERMR}{PERMR}

\title{\textbf{Eclipse 2025.4 — Faults \& Non-Neighbour Connections}\\
       \large 1-to-1 mathematical model with manual equation cross-references}
\author{Library: \texttt{my-axon/libraries/eclipse-faults-nncs}}
\date{Built 2026-05-23 from Eclipse 2025.4 Technical Description and Reference Manual}

\begin{document}
\maketitle

% =====================================================================
\section{Scope and notation}
% =====================================================================

This document gives the formal equations Eclipse uses to compute transmissibilities for:
\begin{itemize}
  \item Block-centred grids (\kw{OLDTRAN}, \kw{OLDTRANR})
  \item Corner-point grids (\kw{NEWTRAN}, the default with \kw{COORD}/\kw{ZCORN})
  \item Half-cell directional permeability (\kw{HALFTRAN}, E100)
  \item Radial grids
  \item Non-Neighbour Connections (NNCs) — both auto-generated and user-defined
\end{itemize}

The \emph{dual-porosity} matrix-fracture transmissibility (\TDref{2.54}) lives in the
sibling document at \texttt{eclipse-dual-porosity/notes/math-model.pdf}.

\paragraph{Symbols.}
\begin{description}
  \item[$\CDARCY$] Darcy unit conversion constant; $0.00112712$ (FIELD), $0.00852702$ (METRIC)
  \item[$\PERMX_i$, $\PERMY_i$, $\PERMZ_i$] permeabilities in cell $i$
  \item[$\RNTG_i$] net-to-gross ratio in cell $i$
  \item[$\TMLT$] composite transmissibility multiplier $= \kw{MULTX}\cdot\kw{MULTY}\cdot\kw{MULTZ}\cdot\MULTFLT\cdot\kw{MULTREGT}\cdot\ldots$
  \item[$\DIPC$] dip correction (for OLDTRAN family)
  \item[$DX_i, DY_i, DZ_i$] in OLDTRAN context: cell side lengths;
                              in NEWTRAN context: vector components from cell centre to face centre
  \item[$A_x, A_y, A_z$] X-, Y-, Z-projections of the mutual interface area between two cells
  \item[$T_i, T_j$] half-cell transmissibilities (used in harmonic mean)
\end{description}

% =====================================================================
\section{Block-centred transmissibility (\kw{OLDTRAN})}
% =====================================================================

Used by default with \kw{DX}/\kw{DY}/\kw{DZ}/\kw{TOPS} input.

\subsection*{X-direction transmissibility}
\begin{equation}
\TRANX_i \;=\; \frac{\CDARCY \cdot \kw{TMLTX}_i \cdot A \cdot \DIPC}{B}
\label{eq:oldtran-x}
\tag{TD 2.1}
\end{equation}
with the interface area
\begin{equation}
A \;=\; \frac{DX_i \cdot DY_j \cdot DZ_j \cdot \RNTG_j \;+\; DX_j \cdot DY_i \cdot DZ_i \cdot \RNTG_i}{DX_i + DX_j}
\label{eq:oldtran-A}
\tag{TD 2.2}
\end{equation}
the flow-resistance sum
\begin{equation}
B \;=\; \frac{\dfrac{DX_i}{\PERMX_i} + \dfrac{DX_j}{\PERMX_j}}{2}
\label{eq:oldtran-B}
\tag{TD 2.3}
\end{equation}
and the dip correction
\begin{equation}
\DIPC \;=\; \frac{\mathrm{DHS}}{\mathrm{DHS} + \mathrm{DVS}}
\label{eq:oldtran-dipc}
\tag{TD 2.4}
\end{equation}
with
\[
\mathrm{DHS} \;=\; \left(\frac{DX_i + DX_j}{2}\right)^{2}
\qquad
\mathrm{DVS} \;=\; \bigl(\mathrm{DEPTH}_j - \mathrm{DEPTH}_i\bigr)^{2}
\]

\subsection*{Y-direction transmissibility}
Analogous to \eqref{eq:oldtran-x} with X $\to$ Y in all subscripts.

\subsection*{Z-direction transmissibility}
\begin{equation}
\TRANZ_i \;=\; \frac{\CDARCY \cdot \kw{TMLTZ}_i \cdot A}{B}
\label{eq:oldtran-z}
\tag{TD 2.5}
\end{equation}
\begin{equation}
A \;=\; \frac{DZ_i \cdot DX_j \cdot DY_j \;+\; DZ_j \cdot DX_i \cdot DY_i}{DZ_i + DZ_j}
\label{eq:oldtran-Az}
\tag{TD 2.6}
\end{equation}
\begin{equation}
B \;=\; \frac{\dfrac{DZ_i}{\PERMZ_i} + \dfrac{DZ_j}{\PERMZ_j}}{2}
\label{eq:oldtran-Bz}
\tag{TD 2.7}
\end{equation}
\note{No dip correction in Z; no net-to-gross factor either ($\RNTG$ is a horizontal-flow correction only).}

% =====================================================================
\section{Alternative block-centred form (\kw{OLDTRANR})}
% =====================================================================

Activated by \kw{OLDTRANR} keyword in the GRID section. Different combination of cell-cross-section
and permeability per cell.

\begin{equation}
\TRANX_i \;=\; \frac{\CDARCY \cdot \kw{TMLTX}_i \cdot \DIPC}{B}
\label{eq:oldtranr-x}
\tag{TD 2.8}
\end{equation}
\begin{equation}
B \;=\; \frac{\dfrac{DX_i}{A_i \cdot \PERMX_i} \;+\; \dfrac{DX_j}{A_j \cdot \PERMX_j}}{2}
\label{eq:oldtranr-B}
\tag{TD 2.9}
\end{equation}
with $A_i = DY_i \cdot DZ_i \cdot \kw{NTG}_i$ and $A_j = DY_j \cdot DZ_j \cdot \kw{NTG}_j$
(per-cell cross-sections).
The dip correction follows \TDref{2.4}/\eqref{eq:oldtran-dipc}.

\subsection*{Z-direction (OLDTRANR)}
\begin{equation}
\TRANZ_i \;=\; \frac{\CDARCY \cdot \kw{TMLTZ}_i}{B}
\label{eq:oldtranr-z}
\tag{TD 2.11}
\end{equation}
\begin{equation}
B \;=\; \frac{\dfrac{DZ_i}{A_i \cdot \PERMZ_i} \;+\; \dfrac{DZ_j}{A_j \cdot \PERMZ_j}}{2}
\label{eq:oldtranr-Bz}
\tag{TD 2.12}
\end{equation}
with $A_i = DX_i \cdot DY_i$ and $A_j = DX_j \cdot DY_j$ (no NTG in vertical).

% =====================================================================
\section{Corner-point transmissibility (\kw{NEWTRAN})}
% =====================================================================

The workhorse formula. Handles both regular face connections and fault-displaced
NNCs uniformly.

\subsection*{X-direction transmissibility}
\begin{equation}
\TRANX_{ij} \;=\; \frac{\CDARCY \cdot \kw{TMLTX}_i}{\dfrac{1}{T_i} + \dfrac{1}{T_j}}
\label{eq:newtran-x}
\tag{TD 2.13}
\end{equation}
The half-cell transmissibility for cell $i$ is
\begin{equation}
T_i \;=\; \PERMX_i \cdot \RNTG_i \cdot \frac{\mathbf{A}\cdot\mathbf{D}_i}{\mathbf{D}_i\cdot\mathbf{D}_i}
\label{eq:newtran-Ti}
\end{equation}
where
\begin{equation}
\mathbf{A}\cdot\mathbf{D}_i \;=\; A_x\,DX_i + A_y\,DY_i + A_z\,DZ_i
\label{eq:newtran-AdotD}
\end{equation}
\begin{equation}
\mathbf{D}_i\cdot\mathbf{D}_i \;=\; DX_i^{\,2} + DY_i^{\,2} + DZ_i^{\,2}
\label{eq:newtran-DdotD}
\end{equation}
with:
\begin{itemize}
  \item $(A_x, A_y, A_z)$ — projections of the mutual interface area onto the global X, Y, Z axes (computed by Eclipse from the corner-point overlap)
  \item $(DX_i, DY_i, DZ_i)$ — components of the vector from cell-$i$ centre to the centre of the face connecting cells $i$ and $j$ (\emph{not} cell side lengths in NEWTRAN context)
\end{itemize}

$T_j$ is computed analogously for cell $j$ (using $\PERMX_j$, $\RNTG_j$, and $\mathbf{D}_j$ — the vector from cell-$j$ centre to the same face centre).

\note{Why this works for NNCs: for an axis-aligned face neighbour pair, $A_y = A_z = 0$ and the formula collapses to the half-cell harmonic mean. For an NNC across a fault, $\mathbf{A}$ encodes the partial overlap geometry (which may have all three components non-zero), and the dot-product structure correctly projects the flow path onto the permeability tensor.}

\subsection*{Y- and Z-direction transmissibility}
Y-transmissibility uses $\PERMY$ in place of $\PERMX$. Z-transmissibility uses $\PERMZ$ \emph{without} the $\RNTG$ factor (vertical flow is between full layer interfaces, no net-to-gross reduction).

% =====================================================================
\section{Half-cell permeability scheme (\kw{HALFTRAN})}
% =====================================================================
Eclipse 100. Allows six effective half-cell permeabilities per cell:
$K_{xil},\,K_{xir},\,K_{yil},\,K_{yir},\,K_{zil},\,K_{zir}$ (left/right of each cell in each direction).

The half-cell transmissibilities are
\begin{equation}
T_{il} \;=\; K_{xil} \cdot \RNTG_i \cdot \frac{\mathbf{A}\cdot\mathbf{D}_i}{\mathbf{D}_i\cdot\mathbf{D}_i}
\label{eq:halftran-Til}
\tag{TD 2.14}
\end{equation}
\begin{equation}
T_{ir} \;=\; K_{xir} \cdot \RNTG_i \cdot \frac{\mathbf{A}\cdot\mathbf{D}_i}{\mathbf{D}_i\cdot\mathbf{D}_i}
\label{eq:halftran-Tir}
\tag{TD 2.15}
\end{equation}
Implemented by storing $\kw{MULTX}_i = K_{xir}/\PERMX_i$ and $\kw{MULTX-}_i = K_{xil}/\PERMX_i$,
with normal $\PERMX$ values entered.

\note{Requires \kw{GRIDOPTS} \kw{YES} in RUNSPEC to enable the negative-direction multipliers (\kw{MULTX-}, \kw{MULTY-}, \kw{MULTZ-}).}

% =====================================================================
\section{Radial transmissibility}
% =====================================================================

For DR/DTHETA/DZ grids:
\begin{equation}
\TRANR_i \;=\; \frac{\CDARCY \cdot \kw{TMLTR}_i \cdot \DIPC}{\dfrac{1}{T_i} + \dfrac{1}{T_j}}
\label{eq:radial-R}
\tag{TD 2.16}
\end{equation}
with pressure-equivalent radii $r_1, r_2, r_3$ giving
\[
T_i \;=\; \frac{\PERMR_i\,\RNTG_i\,D\theta_i\,DZ_i}{F_{1i}}
\qquad
F_{1i} \;=\; \frac{r_1^{\,2}}{r_2^{\,2}-r_1^{\,2}}\,\ln(r_1/r_2) + \tfrac{1}{2}
\]
\[
T_j \;=\; \frac{\PERMR_j\,\RNTG_j\,D\theta_j\,DZ_j}{F_{2j}}
\qquad
F_{2j} \;=\; \frac{r_3^{\,2}}{r_3^{\,2}-r_2^{\,2}}\,\ln(r_3/r_2) - \tfrac{1}{2}
\]
where $r_2$ is the radius of the cell interface and $r_1, r_3$ are pressure-equivalent radii of cells $i$ and $j$.

% =====================================================================
\section{Non-Neighbour Connections}
% =====================================================================

Eclipse auto-generates an NNC between two cells whenever they have a non-zero mutual interface area
but are not Cartesian neighbours. Sources:

\begin{enumerate}
  \item Corner-point fault offset (the canonical case — see below)
  \item Pinch-out bridging (\kw{PINCH})
  \item Local-grid-refinement boundary (\kw{CARFIN}/\kw{RADFIN})
  \item Numerical aquifer connection (\kw{AQUNUM}/\kw{AQUCON})
  \item Dual porosity / dual permeability matrix-fracture coupling (\kw{DUALPORO}/\kw{DUALPERM})
\end{enumerate}

\subsection*{Geometric NNC transmissibility}
For cases 1--3 above, Eclipse computes the NNC transmissibility via the \emph{same} NEWTRAN formula
\eqref{eq:newtran-x}--\eqref{eq:newtran-DdotD}, with the appropriate mutual interface area projections
$(A_x, A_y, A_z)$ for the partial overlap. The formula is uniform across regular faces and NNCs.

\subsection*{Aquifer NNC}
Numerical aquifer transmissibilities use specialised formulas defined by \kw{AQUCON} parameters
(area, length, permeability). Beyond the scope of this document.

\subsection*{Dual-porosity NNC (cross-reference only)}
\begin{equation}
T_{mf} \;=\; \CDARCY \cdot \sigma \cdot K \cdot V
\label{eq:dp-nnc}
\tag{TD 2.54}
\end{equation}
where $\sigma$ is the Kazemi shape factor, $K$ is the matrix permeability, $V$ is the matrix
bulk volume. See sibling library \texttt{eclipse-dual-porosity} for full derivation.

\subsection*{User-defined NNC transmissibility (\kw{NNC} keyword)}
The user enters the transmissibility directly:
\begin{equation}
T_{ij}^{\text{user}} \;=\; \kw{TRAN} \quad \text{(item 7 of the \kw{NNC} record)}
\label{eq:nnc-user}
\end{equation}

If the user also supplies an area $A_{ij}$ (item 16) and linking permeability $K^{\text{link}}_{ij}$
(item 17) — both required for non-Darcy flow (\kw{VDFLOW}) — Eclipse back-computes the half-cell distances:
\begin{equation}
DX_i + DX_j \;=\; \frac{2\,A_{ij}\,K^{\text{link}}_{ij}\,c_{\text{conv}}}{T_{ij}^{\text{user}}}
\label{eq:nnc-backcompute}
\end{equation}
where $c_{\text{conv}}$ is a unit conversion factor. Implementation split: \(DX_i = DX_j = (DX_i + DX_j)/2\).

\subsection*{Combined multiplier semantics}
For any connection between cells $i$ and $j$, the final transmissibility is the product:
\begin{equation}
T_{\text{final}}\;=\; T_{\text{base}} \;\cdot\; M_{i,j}^{\,(\text{cell})} \;\cdot\; M^{\,\text{fault}} \;\cdot\; M^{\,\text{region}} \;\cdot\; M^{\,\text{EDITNNC}}
\label{eq:Tfinal}
\end{equation}
where
\begin{itemize}
  \item $T_{\text{base}}$ is the geometric Tx from NEWTRAN \eqref{eq:newtran-x} or OLDTRAN \eqref{eq:oldtran-x};
  \item $M_{i,j}^{\,(\text{cell})}$ is the product of per-cell \kw{MULTX/Y/Z} (and \kw{MULTX-/Y-/Z-} where applicable);
  \item $M^{\,\text{fault}}$ is the \kw{MULTFLT} multiplier for the named fault face containing this connection (1 if no fault);
  \item $M^{\,\text{region}}$ is the \kw{MULTREGT} multiplier between the cells' \kw{MULTNUM} regions (1 if same region);
  \item $M^{\,\text{EDITNNC}}$ is the \kw{EDITNNC} multiplier (1 if not edited; replaced by a constant if \kw{EDITNNCR} is used).
\end{itemize}

% =====================================================================
\section{Dip correction $\DIPC$}
% =====================================================================
For block-centred OLDTRAN/OLDTRANR:
\begin{equation}
\DIPC \;=\; \frac{\mathrm{DHS}}{\mathrm{DHS}+\mathrm{DVS}}
\quad,\quad
\mathrm{DHS} = \left(\tfrac{DX_i + DX_j}{2}\right)^{2}
\quad,\quad
\mathrm{DVS} = \bigl(\mathrm{DEPTH}_j - \mathrm{DEPTH}_i\bigr)^{2}
\label{eq:dipc-full}
\tag{TD 2.4/2.10}
\end{equation}
\textbf{Limit cases}: horizontal layer $\Rightarrow$ DVS = 0 $\Rightarrow$ $\DIPC = 1$. Highly dipping
$\Rightarrow$ $\DIPC < 1$, reducing effective Tx. Applied only to X and Y (not Z).
NEWTRAN does not require an explicit $\DIPC$ — the corner-point geometry already encodes dip in
$\mathbf{D}_i, \mathbf{D}_j$.

% =====================================================================
\section{Equation \(\to\) keyword cross-reference}
% =====================================================================

\begin{center}\small
\begin{tabular}{lll}
\toprule
Equation & Concept & Controlling keyword(s) \\
\midrule
\eqref{eq:oldtran-x}--\eqref{eq:oldtran-dipc} & Block-centred X Tx (OLDTRAN)             & \kw{OLDTRAN}, \kw{DX}, \kw{DY}, \kw{DZ}, \kw{PERMX}, \kw{MULTX} \\
\eqref{eq:oldtran-z}--\eqref{eq:oldtran-Bz}   & Block-centred Z Tx                       & \kw{OLDTRAN}, \kw{DZ}, \kw{PERMZ}, \kw{MULTZ}                  \\
\eqref{eq:oldtranr-x}--\eqref{eq:oldtranr-Bz} & Alternative block-centred                & \kw{OLDTRANR}                                                  \\
\eqref{eq:newtran-x}--\eqref{eq:newtran-DdotD} & Corner-point Tx (NEWTRAN, NNC formula)  & \kw{NEWTRAN}, \kw{COORD}, \kw{ZCORN}                            \\
\eqref{eq:halftran-Til}--\eqref{eq:halftran-Tir} & Half-cell perm scheme                 & \kw{HALFTRAN}, \kw{MULTX-}, \kw{MULTY-}, \kw{MULTZ-}, \kw{GRIDOPTS} \\
\eqref{eq:radial-R}                            & Radial Tx                              & \kw{DR}, \kw{DTHETA}, \kw{INRAD}                                \\
\eqref{eq:nnc-user}--\eqref{eq:nnc-backcompute} & User-defined NNC                       & \kw{NNC}, \kw{NNCGEN}                                            \\
\eqref{eq:dp-nnc}                              & DP/DK matrix-fracture NNC               & \kw{DUALPORO}/\kw{DUALPERM}, \kw{SIGMA}, \kw{SIGMAV}             \\
\eqref{eq:Tfinal}                              & Combined multiplier product            & \kw{MULTX}, \kw{MULTFLT}, \kw{MULTREGT}, \kw{EDITNNC}, \kw{EDITNNCR} \\
\eqref{eq:dipc-full}                           & Dip correction (OLDTRAN family only)    & implicit; not user-controlled                                   \\
\bottomrule
\end{tabular}
\end{center}

% =====================================================================
\section{Default \kw{CDARCY} values}
% =====================================================================

\begin{center}\small
\begin{tabular}{lll}
\toprule
Unit system & E100 value & E300 value \\
\midrule
METRIC  & 0.008527 & 0.00852702 \\
FIELD   & 0.001127 & 0.00112712 \\
LAB     & 3.6      & 3.6        \\
PVT-M   & 0.00864  & 0.00864    \\
\bottomrule
\end{tabular}
\end{center}

These conversion constants give transmissibilities in units of cP$\cdot$reservoir-volume / (day $\cdot$ pressure) — specifically:
$\text{cP}\cdot\text{rm}^{3}/\text{day}/\text{bars}$ (METRIC),
$\text{cP}\cdot\text{rb}/\text{day}/\text{psi}$ (FIELD),
$\text{cP}\cdot\text{rcc}/\text{hr}/\text{atm}$ (LAB),
$\text{cP}\cdot\text{rm}^{3}/\text{day}/\text{atm}$ (PVT-M).

% =====================================================================
\section{NNC drop / inline thresholds}
% =====================================================================
\begin{itemize}
  \item NNCs with $T < 10^{-6}$ are silently dropped from the connection list (RM p.1468).
  \item NNCs with $T < 10^{-20}$ produce a warning and are dropped.
  \item NNCs whose path is geometrically equivalent to a regular face connection (e.g.\ entire intervening layer inactive) are \emph{inlined} — converted back to regular face connections internally. Visible only via the \kw{ALLNNC} mnemonic in \kw{RPTGRID}.
\end{itemize}

% =====================================================================
\section{Composition with the DP/DK system}
% =====================================================================

For a fully coupled DP+faulted run, the Jacobian contains connections from:
\begin{enumerate}
  \item Standard NEWTRAN/OLDTRAN face-neighbour connections (fracture-fracture and, under DK, matrix-matrix)
  \item Auto-generated geometric NNCs from corner-point fault offsets — same NEWTRAN formula with partial overlap area
  \item Pinch-out NNCs from \kw{PINCH}
  \item Dual-porosity $\sigma$-NNCs (separate Kazemi formula, \eqref{eq:dp-nnc})
  \item Block-to-block NNCs from \kw{BTOBALFA} (standard spatial Tx with area multiplier)
\end{enumerate}

\kw{MULTFLT} applies to (1), (2), (3), and (5). It does \emph{not} apply to (4) (DP $\sigma$-NNCs);
modify those via \kw{SIGMAV} or \kw{MULTSIGV}.

% =====================================================================
\appendix
\section{Equation index (by TD equation number)}
% =====================================================================
\begin{center}\small
\begin{tabular}{ll}
\toprule
TD Eq. & Label in this document \\
\midrule
2.1  & \eqref{eq:oldtran-x}  \\
2.2  & \eqref{eq:oldtran-A}  \\
2.3  & \eqref{eq:oldtran-B}  \\
2.4  & \eqref{eq:oldtran-dipc} \\
2.5  & \eqref{eq:oldtran-z}  \\
2.6  & \eqref{eq:oldtran-Az} \\
2.7  & \eqref{eq:oldtran-Bz} \\
2.8  & \eqref{eq:oldtranr-x} \\
2.9  & \eqref{eq:oldtranr-B} \\
2.10 & \eqref{eq:dipc-full}  \\
2.11 & \eqref{eq:oldtranr-z} \\
2.12 & \eqref{eq:oldtranr-Bz}\\
2.13 & \eqref{eq:newtran-x}  \\
2.14 & \eqref{eq:halftran-Til}\\
2.15 & \eqref{eq:halftran-Tir}\\
2.16 & \eqref{eq:radial-R}   \\
\bottomrule
\end{tabular}
\end{center}

\bigskip\noindent\emph{End of mathematical model.}

\end{document}
